\section{OPTICAL PLATFORMS AND INTEGRATION ARCHITECTURES}
\label{sec:optical_modality_families}
\label{sec:integration_architectures}

In O-ISAC, the physical role of optics determines both where information
propagates and where the environment is observed. Optics may carry the user
link, illuminate a target, distribute a wireless carrier, or provide a coherent
reference, and each role exposes a different set of impairments. The integration
architecture then reveals which component, signal, or resource is shared by
communication and sensing. We therefore examine the evidence along the axes of
physical platform and coupling mechanism.

Table~\ref{tab:modality_map} reports the mutually exclusive family counts and
corpus shares; photonics-assisted terahertz and fiber together account for 125
studies (60.7\% of the corpus).

\begin{table}[t]
\caption{Canonical physical organization of the 206-study corpus}
\label{tab:modality_map}
\centering
\footnotesize
\setlength{\tabcolsep}{4pt}
\renewcommand{\arraystretch}{1.06}
\begin{tabularx}{\columnwidth}{@{}>{\raggedright\arraybackslash}X r@{}}
\toprule
Platform family & Studies, $n$ (\%) \\
\midrule
Photonics-assisted THz & 69 (33.5\%) \\
Fiber & 56 (27.2\%) \\
VLC/LiFi & 38 (18.4\%) \\
Free-space optical & 31 (15.0\%) \\
Hybrid optical & 9 (4.4\%) \\
Other optical & 3 (1.5\%) \\
\midrule
\textbf{Total} & \textbf{206 (100.0\%)} \\
\bottomrule
\end{tabularx}
\vspace{1pt}
\parbox{\columnwidth}{\footnotesize\emph{Note.} Each study is counted once under
its dominant platform family. Media and integration mechanisms can overlap
within the represented systems. The counts summarize how the reviewed evidence
is distributed across the six platform families.}
\end{table}

The family profiles below provide the corresponding physical paths,
observables, coupling mechanisms, active constraints, transfer conditions, and
evidence anchors. Each study belongs to one platform family, while its coupling
mechanisms may overlap. Transmitters supplied through fiber enter the terahertz
family when photonic generation and conversion create the wireless carrier
directed at the target. Fiber remains the organizing platform when it transports
signals to downstream radio-frequency identification (RFID), radio, or radar
sensing. Hybrid membership begins when interaction between physical segments
shapes the result. This organization locates the target path and guides the
synthesis that follows.

\subsection{Physical Platform Families}

\subsubsection{Photonics Assisted Terahertz Systems}

Photonics-assisted terahertz is the largest family, comprising 69 of the 206
studies (33.5\%). The optical chain generates, distributes, or processes a
coherent carrier before conversion to a wireless terahertz or millimeter-wave
target path, from which delay, range, Doppler, angle, imaging, and channel- or
target-response observables are obtained. Communication and sensing recur as
hardware-, reference-, waveform-, beam-, or resource-coupled functions
\cite{OISAC_SCR01054,OISAC_SCR00036,OISAC_SCR00083}. This upstream sharing ties
both functions to the same optical conditions. Laser coherence and linewidth
govern reference stability, while component bandwidth, conversion fidelity,
and efficiency shape the bandwidth, phase structure, and power delivered to
the antenna \cite{OISAC_SCR00988,OISAC_SCR01004}. At the radio interface,
antenna aperture, gain, and alignment govern launch and collection, molecular
absorption shapes propagation, and target scattering forms the returning echo
in radar configurations \cite{OISAC_SCR00989,OISAC_SCR00992}.

Where the functions meet varies by architecture, with some systems dividing a
multicarrier frame between sensing and communication and others placing the
joint decision in a phased array controlled through photonics
\cite{OISAC_SCR01054,OISAC_SCR00036}. Wavelength tuning and multibeam reception
move coordination to the source or the receiving front end
\cite{OISAC_SCR00185,OISAC_SCR00396}. At the waveform level, one multichannel
design combines linear frequency modulation (LFM) with communication symbols
across coherent carriers \cite{OISAC_SCR00988}. Other designs add source
agility through frequency hopping or organize the signal over delay and Doppler
through orthogonal multiplexing \cite{OISAC_SCR00992,OISAC_SCR00973}.
Information can also close the loop, as burst SNR estimates guide both
probabilistic shaping and ranging, while a simulated radar waveform carries
communication data \cite{OISAC_SCR00989,OISAC_SCR00882}. Together, these
designs show how the joint decision can move from hardware to waveform and to
the information used by both tasks.

\FloatBarrier

The location of that decision also determines how the measurement is formed. A
returned echo may be remodulated onto light and carried back over fiber for
centralized processing, where photonic dechirping expresses target delay at a
lower frequency \cite{OISAC_SCR00966,OISAC_SCR00992}. Along this path, phase
noise associated with laser linewidth degrades the recovered signal, while
delay mismatch between optical and electrical paths shifts the ranging
observable \cite{OISAC_SCR00966}.
Laboratory experiments reveal waveform, conversion, and beamforming behavior
under controlled conditions \cite{OISAC_SCR00988,OISAC_SCR00036}, whereas a
270~m outdoor link adds wireless propagation and operation of the complete
chain \cite{OISAC_SCR01054}. Meaningful numerical transfer therefore requires
aligned optical and wireless stages, aperture and beam conditions, receiver
processing, and validation setting.

\subsubsection{Fiber O-ISAC}

Fiber comprises 56 studies (27.2\%) and follows two physical paths. The guided
link may form the sensing medium, with disturbances observed through
backscatter, carrier phase, modal coupling, or an inserted sensor
\cite{OISAC_SCR00957,OISAC_SCR00432,OISAC_SCR00356,OISAC_SCR00679}. Fiber may
instead transport signals to a downstream RFID, radio, or radar segment, so the
target interaction occurs after optical distribution
\cite{OISAC_SCR00496,OISAC_SCR00631,OISAC_SCR00719,OISAC_SCR00983,OISAC_SCR01049}.
The corresponding observations include distributed events, vibration, strain,
and environmental change in the guided path, together with downstream tag
state, range, localization, and imaging
\cite{OISAC_SCR00957,OISAC_SCR00631,OISAC_SCR00496}. The measurement boundary
therefore lies either within the guided path or across the complete optical and
wireless chain.

Coupling recurs through shared hardware, wavelengths, waveforms, links, or
resource decisions. Integration can reuse the communication signal itself
\cite{OISAC_SCR00647,OISAC_SCR00038,OISAC_SCR00356,OISAC_SCR00432,OISAC_SCR00015}
or place distinct data and sensing signals on shared infrastructure
\cite{OISAC_SCR00013,OISAC_SCR00344,OISAC_SCR00957,OISAC_SCR00218,OISAC_SCR00679}.
The first form ties sensing to the symbols, carrier, phase, or supported modes,
whereas the second ties both functions to the same fiber and receiver
conditions. Guided sensing is constrained by dynamic range, leakage,
reflection, nonlinearity, reference stability, topology, and traffic, whereas
transported sensing additionally depends on conversion, timing, remote
calibration, and wireless geometry. Evidence from laboratory links, live
access networks, and a deployed submarine route shows why numerical transfer
must match the interrogation or transport role, topology, launch and receiver
conditions, conversion and clocking, target path, traffic, and validation
setting
\cite{OISAC_SCR00647,OISAC_SCR00038,OISAC_SCR00218,OISAC_SCR00957,
OISAC_SCR00631,OISAC_SCR00496,OISAC_SCR01049}.

\subsubsection{VLC/LiFi Systems}

VLC/LiFi comprises 38 studies (18.4\%) and places communication and sensing
within the same illuminated scene. The grouping preserves the distinct link-
and networking-level meanings of VLC and LiFi. Direct or reflected room and
device paths through emitters, cameras, and photodiodes support position, pose,
distance, imaging, gesture, recognition, and physiological sensing
\cite{OISAC_SCR00910,OISAC_SCR01034,OISAC_SCR00196}. Cameras and photodiodes
expose different balances between spatial detail and reception speed, while
nonnegativity, dimming, flicker, source linearity, ambient light, field of view,
and exposure constrain what can be communicated and observed
\cite{OISAC_SCR00273,OISAC_SCR00925,OISAC_SCR00910,OISAC_SCR01036}.
Coupling recurs through hardware, waveform, geometry, resource decisions, and
the application loop. It may arise through a shared observation
\cite{OISAC_SCR01034,OISAC_SCR00836}, through geometry and resource decisions
\cite{OISAC_SCR00901,OISAC_SCR00905,OISAC_SCR00407,OISAC_SCR00557,OISAC_SCR00915},
or through integrated emitters and detectors
\cite{OISAC_SCR00943,OISAC_SCR00873,OISAC_SCR00001}.

Illumination is the common comparison boundary across these designs. Geometry
and multipath shape localization, reflected observations support gesture and
recognition, and learned models extend the same optical field to health sensing
tasks
\cite{OISAC_SCR00910,OISAC_SCR00407,OISAC_SCR00557,OISAC_SCR01036,
OISAC_SCR00376,OISAC_SCR01047}. Analytical and
simulated studies isolate geometry and resource choices, whereas prototypes and
field experiments add lighting, receiver, and control conditions
\cite{OISAC_SCR00901,OISAC_SCR00407,OISAC_SCR00915,OISAC_SCR00910,
OISAC_SCR00836,OISAC_SCR00873}. Numerical transfer requires matched task
definitions, layout, orientation, illumination, detector, processing, and
validation conditions.

\subsubsection{Free Space Optical Systems}

Free-space optical systems comprise 31 studies (15.0\%) and use a directional
link, optionally combined with target illumination and return, over direct or
reflected aperture and beam paths. These paths support ranging, localization,
imaging, tracking, and atmospheric or channel observation
\cite{OISAC_SCR00011,OISAC_SCR01029,OISAC_SCR00730}. Geometry is therefore part
of both the link and the measurement. Coupling recurs through the link,
waveform, aperture or beam, and spatial or temporal resources, while pointing,
divergence, aperture, reflectivity, loss, turbulence, and motion constrain the
resulting evidence. Direct, reflected, scanned, and steered paths illustrate
these relations \cite{OISAC_SCR00011,OISAC_SCR00012,OISAC_SCR00506,OISAC_SCR00941}.
Retroreflection and multicarrier reception connect returned power to range and
position \cite{OISAC_SCR01029,OISAC_SCR00730}. Structured waveforms and
steering apertures place coupling in the signal or beam
\cite{OISAC_SCR00256,OISAC_SCR00903,OISAC_SCR00904,OISAC_SCR00594}, while
atmospheric estimates can guide beam shaping or link selection
\cite{OISAC_SCR01085,OISAC_SCR00456}.

Across these forms, beam geometry determines both spatial selectivity and the
power collected at the receiver. Range, angle, and image features
therefore remain tied to calibration, aperture, pointing, and motion
\cite{OISAC_SCR00011,OISAC_SCR00012,OISAC_SCR00506,OISAC_SCR01029,
OISAC_SCR00730}. Controlled links reveal waveform and estimator behavior,
whereas operational ranging adds atmospheric variation and continuous field
observations \cite{OISAC_SCR01085,OISAC_SCR00730}. Numerical transfer requires
matched geometry, calibration, target properties, weather, motion, and
validation conditions. Hybrid systems extend that comparison across interfaces
between physical segments.

\subsubsection{Hybrid Optical Systems}

Hybrid optical systems comprise nine studies (4.4\%) in which consequential
interaction occurs among optical, fiber, free-space, or radio segments and the
perturbation and observation may straddle a conversion boundary. Their
observables include range, vibration, and target or network state within
individual segments, while coupling arises through a central reference and
through links, resources, control, or applications shared across segments.
Conversion loss, segment power, synchronization, calibration, and timing are
therefore active constraints
\cite{OISAC_SCR00210,OISAC_SCR00366,OISAC_SCR00277}. Numerical transfer must
trace every segment, conversion, clock, and measurement plane and validate each
communication and sensing function.

Unlike fiber used only for downstream transport, hybrid membership requires
conversion or cross-segment feedback to shape the result. This occurs when
fiber-to-wireless interaction is consequential
\cite{OISAC_SCR00210,OISAC_SCR00366}, when a shared photonic source spans fiber,
free space, and millimeter-wave stages \cite{OISAC_SCR00944}, or when sensing in
one segment guides service in another \cite{OISAC_SCR01011,OISAC_SCR00277}.
Other designs establish a common signal origin through shared time modes or a
frequency-agile photonic waveform
\cite{OISAC_SCR00091,OISAC_SCR01229}.

\subsubsection{Other Optical Systems}

The three other optical studies (1.5\%) use distinct optical-control or
coherent-carrier mechanisms: a microcomb-generated microwave field, a
solar-blind ultraviolet path, and a quantum-optical path. Their observables are,
respectively, a field pattern, particulate or smoke response, and phase
\cite{OISAC_SCR00268,OISAC_SCR00379,OISAC_SCR00906}. Coupling is specific to
specialized hardware, reference, and application choices, and each case is
bounded by its own propagation, target, receiver, and measurement models.
Numerical results must therefore be interpreted within the native mechanism and
measurement model of each system. Together, the hybrid and other optical cases
complete the physical map before the analysis turns to coupling mechanisms
shared across families.

\subsection{Integration Architectures and Resource Coupling}

Across these physical families, integration architecture locates where
communication and sensing meet along the target path. The coupling may enter
through shared hardware, a common carrier or waveform, coordinated resources,
a propagation link, or a service decision. Several mechanisms can coexist in
one system. Each label identifies one coupling location while preserving the
rest of the chain for separate interpretation.

Coverage is therefore multilabel. Shared hardware appeared in 117 studies and
shared optical carriers in 49. Shared waveforms appeared in 113 studies,
coordinated resource allocation in 118, and shared link or channel use in 87.
Joint design or optimization appeared in 72 studies, and application coupling
in 46. The normalized map covers all 206 studies and retains three boundary
cases as mixed.

% FUTURE FIGURE 5 --- MULTILABEL INTEGRATION AND RESOURCE COUPLING MAP
% Do not generate in this revision.
% Stable label: fig:integration_map
% Blueprint ID: FIG-OISAC-04
% Placement: immediately after a compact interpretive lead-in and before the
% mechanism subsections.
% Activation rule: once the figure is live, replace the current seven-count
% prose inventory with the two-sentence lead-in below. Do not repeat the counts
% in the surrounding prose.
% Data authority: Phase F s2_integration_mechanisms.csv and its locked
% normalization audit; denominator 206, multilabel.
% Layout: a generic left-to-right system chain with parallel communication and
% sensing lanes: source/front end -> carrier/modulation -> waveform/frame ->
% resource controller -> propagation link/channel and target -> receiver/DSP ->
% service/application. Seven equal coupling bands connect the lanes at their
% usual locations: shared hardware 117; optical carrier 49; waveform 113;
% resource allocation 118; link or channel 87; joint design or optimization 72;
% application scenario 46. Retain three mixed boundary cases as a separate
% audit callout. Node order shows location, not frequency or maturity; badge size
% remains constant and counts must not be summed.
% Future lead-in: ``Figure~\ref{fig:integration_map} locates the seven recurring
% integration mechanisms along a generic O-ISAC signal and service path. The
% badges report overlapping coverage rather than an exclusive taxonomy.''
% Exact future caption: ``Locations of communication and sensing coupling in
% the 206-study O-ISAC corpus. Shared resource allocation (118 studies),
% hardware (117), waveform (113), link or channel (87), joint design or
% optimization (72), optical carrier (49), and application scenario (46) are
% shown at their typical locations; three mixed boundary cases remain separate.
% Coding is multilabel, so counts must not be summed and do not indicate
% integration strength or maturity.''

\subsubsection{Shared Hardware, Carrier, and Link}

Hardware, carrier, and link sharing describe different physical relations.
Reusing a component makes its bandwidth, coherence, dynamic range, or alignment
conditions relevant to both functions. A common optical carrier adds a phase or
wavelength relation, while a common link adds shared propagation and calibration
conditions. Together, these relations show which design choice is common and
which constraint follows it.

The physical path gives reuse its engineering meaning. In fiber, traffic and
sensing share the guided structure, with a resonance shift, phase, or
backscatter conveying the sensed change
\cite{OISAC_SCR00002,OISAC_SCR00013}. A retroreflective free space
path supports information delivery together with ranging or localization
\cite{OISAC_SCR00011}. Arrays controlled through photonics instead share
optical generation and beam control while the target path remains wireless
\cite{OISAC_SCR00036}. Comparison therefore follows the component role, carrier
relation, propagation path, and calibration.

\subsubsection{Shared Waveforms}

Waveform sharing begins when one signal structure contributes to both decoding
and sensing. The shared element may be the complete waveform or a selected
pilot, chirp, code, polarization state, or spatial pattern. This common design
links payload efficiency, sensing observability, and receiver processing.
Structural sharing records that relationship. Concurrent performance adds the
requirement that both outcomes be observed at a matched operating point, a
distinction carried into Section~\ref{sec:metrics}.

The evidence illustrates several forms. Pilots inserted for channel estimation
can also provide a sensing reference \cite{OISAC_SCR00007}. One photonic
terahertz prototype generates a waveform whose frequency steps provide the
radar dimension and whose amplitude carries data. Its experimental evidence is
strongest for waveform generation and communication response
\cite{OISAC_SCR00160}. A visible light design combines temporal data with
spatially structured illumination and uses localization to adapt beamforming
\cite{OISAC_SCR00196}. The shared structure is similar, but its physical and
receiver implementation remains platform specific.

\subsubsection{Resource Allocation and Joint Design}

Resource coupling coordinates temporal, spectral, power, spatial, or processing
degrees of freedom. An allocation can change payload opportunity, link margin,
sensing observability, geometry, or latency. Joint design connects the two
functions more directly by choosing a variable, objective, or constraint in
view of both outcomes. The resource, allocation rule, and operating conditions
therefore define the resulting evidence.

A visible light V2X camera model makes exposure time a shared decision variable.
Its coexistence result is analytical and numerical, while laboratory
communication and sensing gains were obtained in separate configurations
\cite{OISAC_SCR00273}. In one reported terahertz configuration, a rapidly swept
photonic source and common waveform plan support simultaneous data transmission
and ranging \cite{OISAC_SCR00083}. These cases distinguish a modeled optimum
from a measured concurrent configuration. Section~\ref{sec:tradeoffs} therefore
reads each relationship within its operating region and evidence form.

\subsubsection{Coupling at the Application Level and Boundary Cases}

Application coupling arises when sensing informs beam or link selection,
trajectory planning, blockage response, or network control. The shared element
is then the system action, even when the signal paths remain separate. A near
infrared wearable combines communication and physiological sensing on a common
material platform while retaining separate detector roles
\cite{OISAC_SCR00001}. A hybrid last mile design places free space optical
delivery and fiber Bragg grating monitoring within one system model
\cite{OISAC_SCR00277}. Their coordination lies in hardware and service context.

Three studies cross temporal or physical segment boundaries. One performs
localization and imaging before data delivery on a shared optical source and
path with separate receivers \cite{OISAC_SCR00012}. Two use sensed free space
optical channel state to adapt trajectory and resource allocation or cell
placement and user association across optical backhaul and radio access
\cite{OISAC_SCR00917,OISAC_SCR00922}. The mixed label preserves this path and
keeps each result tied to the segment, function, and operating stage that
produces it.

\subsection{Interpretation Across Platforms}

Read together, the platform and architecture views form one mechanism map. The
platform establishes the physical transfer conditions, while the architecture
identifies the shared choice. In fiber, waveform sharing inherits guided
propagation and reflections, together with launch and receiver conditions
\cite{OISAC_SCR00038,OISAC_SCR00647}. In free space optical systems, the same
mechanism is shaped by aperture, geometry, and alignment, whereas terahertz
systems assisted by photonics also introduce conversion from optics to radio
\cite{OISAC_SCR00011,OISAC_SCR00036}. Application coupling can instead
coordinate separate optical and radio paths through a common service decision
\cite{OISAC_SCR00917}.

Across these mechanisms, sharing concentrates design around a common component,
reference, or resource. That choice can couple the functions through power
balance and interference, while synchronization and receiver dynamic range
constrain their joint operation
\cite{OISAC_SCR00013,OISAC_SCR00366}. Section~\ref{sec:metrics} therefore
examines which communication and sensing outcomes change together at a matched
operating point and which conditions explain that relationship.
